% ============================================================
%  第七部分：诚实的边界与修正判断
% ============================================================

\section{诚实的边界：反方证据与修正判断}
\label{sec:boundaries}

前面的章节建立了方法、验证与工程体系。作为教程的收束，本章必须完成一件很多教程回避的工作：\textbf{系统性地呈现反方证据}——Regime 建模在什么条件下失效、被高估、甚至有害。这不是谦虚的姿态，而是第一部分“信息论下界”与第三部分“验证体系”的必然延伸：一个承认滞后下界的框架，也必须承认自身价值的边界。

\subsection{五条反方证据}

\subsubsection{证据一：拥挤与衰减}

任何被广泛使用的信号都会衰减，Regime 信号不例外。当“状态切换”成为行业共识时，切换时点本身会成为交易对象：防御状态的确认信号被抢先交易，切换的一致性放大流动性冲击——这本质上是第一类非平稳（结构演化）在 Regime 系统自身的投影。更一般地，Harvey et al.~\cite{harvey2016} 关于金融研究中数据窥探的系统性证据提醒我们：\emph{从大量试验中筛选出的“有效模式”，其样本外价值通常远低于样本内读数}。Regime 系统若经过大量窗口/状态数/阈值的筛选，其“有效”结论必须按同一标准折价。

\subsubsection{证据二：样本内解释力 \texorpdfstring{$\ne$}{不等于} 预测力}

HMM 能对历史给出漂亮的状态划分，这几乎不构成证据：任何有足够状态数的模型都能“解释”历史。真正的判据是第三部分的三角验证（模型自洽、结构稳定、经济价值），且经济价值必须在\emph{严格因果}的条件下度量。大量研究报告的 Regime 回测恰恰在这一点上失守（平滑泄漏、全样本参数、缺延迟）——这不是个别错误，而是跨文献的系统性方法论缺陷。

\subsubsection{证据三：状态数的不可辨识}

第二部分已论证：$K$ 的选择是决策分辨率的选择，不是对“真实状态数”的逼近；信息准则常在 $K=2\sim3$ 后改善微弱，而经济可解释性与切换成本往往给出更硬的约束。任何声称“市场确实有 5 个状态”的论述都是过度解读；正确的表述是“在可承受的切换成本下，5 个状态提供了足够的分辨率”。

\subsubsection{证据四：结构演化的持续侵蚀}

第一章把非平稳分成三层：参数漂移、状态切换、结构演化。前两层是 Regime 建模的作用域，第三层不是——它没有可重演的样本。市场微观结构变革、参与者结构变化（被动化、量化渗透率上升）、监管框架调整，会缓慢改变每个状态的统计画像本身。一个在 2016--2020 年训练的“高波动状态”画像，在 2024 年可能已经指向不同的资产行为。这不是模型的失败，而是作用域的边界；应对方式只能是监控（画像漂移检测）与定期重校准，而不是期望一个模型一劳永逸。

\subsubsection{证据五：诊断能力的脆弱性}

无标签世界里的评估依赖间接证据，而间接证据可以被“表演”：优化到漂亮的画像（过拟合状态边界）、在历史上选出最平滑的窗口组合（选择偏差）、用泄漏版回测展示优越性（有意或无意的）。第三部分的检查清单之所以必要，正是因为\textbf{这些表演在表面上与真本事无法区分}——只有通过对照实验（合规版 vs 泄漏版）、扰动检验与多重检验折扣，才能让两者重新可分。

\subsection{修正后的判断：这个系统到底值得做吗}

把全部证据合起来，得到一个修正后的价值定位：

\begin{enumerate}[itemsep=3pt]
\item \textbf{Regime 系统不是预测机器}。它不承诺“更早看到未来”，第一部分的下界论证排除了这种可能性。任何以“提高预测精度”为卖点推销 Regime 建模的论述都应当被怀疑。
\item \textbf{它的真实价值有三层}：
\begin{itemize}[itemsep=2pt]
\item \emph{风险管理的纪律化}：把“什么时候该缩小暴露”从情绪判断变成程序判断（Schwert 的波动率时代、Scherer 的传染证据都指向“防御的价值主要在风险端”）；
\item \emph{策略调度的结构化}：因子权重、杠杆上限、止损参数随状态条件化（第~\ref{sec:strategy}~节的证据：条件化配置显著改善风险调整表现）；
\item \emph{沟通与复盘的语言}：把“最近很难做”翻译为“转折点密度上升、趋势结构恶化”——可量化、可跟踪、可追责。
\end{itemize}
\item \textbf{期望值的诚实锚点}：现有文献的一致信号是\emph{风险特征改善}（波动率、回撤、Calmar）与\emph{实施稳健性}（换手、延迟鲁棒性），而不是收益中枢的抬升。以此为期望值做投入产出评估；若团队的目标是“找到更强的预测信号”，Regime 建模是错误的工具。
\item \textbf{明确的适用边界}：
\begin{itemize}[itemsep=2pt]
\item 适合：资产配置层的姿态与风险预算、多策略资金分配、风控熔断、研究框架与复盘语言；
\item 慎用：作为独立 alpha 信号、日内级别的高频择时、缺少成本约束下的高频状态切换；
\item 禁用场景：以平滑概率做实时决策、以全样本参数做历史回测、以单次样本内结果做资本承诺。
\end{itemize}
\end{enumerate}

\subsection{生产纪律清单（十二条）}

作为全教程的操作性收束，以下清单可直接用于团队审查：

\begin{enumerate}[itemsep=2.5pt]
\item 状态判定只使用\emph{滤波}概率；平滑与离线方法只用于研究与复盘；
\item 参数估计使用滚动/扩展窗口逐步重估，禁止全样本参数进入历史回测；
\item 每一期报告\emph{状态画像}（统计签名）而非状态编号；
\item 状态数 $K$ 的选择同时满足经济可解释性与“肘部法则”，并评估切换成本；
\item 跨窗口追踪状态身份一致性（模板匹配/画像指纹），匹配率骤降即人工复核；
\item 泄漏审计（合规版 vs 泄漏版对照实验）纳入常规回归测试；
\item 稳定性检验（子样本重估、种子扰动、块 bootstrap 区间）常态化输出；
\item 姿态响应不对称：进入防御快、恢复进攻慢；
\item 状态概率连续映射到仓位，避免二值开关；
\item 全部回测计入交易成本与 1--2 期执行延迟；
\item 对多重检验保持折扣意识（显著性的门槛应远高于名义 $t=2$）；
\item 每年做一次结构复核：检测状态画像漂移，决定是否重新校准或降级模型。
\end{enumerate}

\subsection{最后的主句}

\keynote{Regime 系统最诚实的使用方式，是把它当作“风险管理的纪律装置”与“策略调度的语法”，而不是一台预测机器。它能显著改善你“保持姿态”的一致性，但无法让你“看得更远”——它全部的边界与全部的价值，都来自同一个事实。}